%%%%%%%% ICML 2026 SUBMISSION ON FISHER-GUIDED SPARSITY-AWARE MODEL MERGING %%%%%%%%%%%%%%%%%

\documentclass{article}

% Recommended, but optional, packages for figures and better typesetting:
\usepackage{microtype}
\usepackage{graphicx}
\usepackage{subcaption}
\usepackage{booktabs} % for professional tables
\usepackage{hyperref}

% Attempt to make hyperref and algorithmic work together better:
\newcommand{\theHalgorithm}{\arabic{algorithm}}

% Use the following line for the initial blind version submitted for review:
\usepackage{icml2026}

\usepackage{amsmath}
\usepackage{amssymb}
\usepackage{mathtools}
\usepackage{amsthm}
\usepackage{algorithm}
\usepackage{algorithmic}

% if you use cleveref..
\usepackage[capitalize,noabbrev]{cleveref}

%%%%%%%%%%%%%%%%%%%%%%%%%%%%%%%%
% THEOREMS
%%%%%%%%%%%%%%%%%%%%%%%%%%%%%%%%
\theoremstyle{plain}
\newtheorem{theorem}{Theorem}[section]
\newtheorem{proposition}[theorem]{Proposition}
\newtheorem{lemma}[theorem]{Lemma}
\newtheorem{corollary}[theorem]{Corollary}
\theoremstyle{definition}
\newtheorem{definition}[theorem]{Definition}
\newtheorem{assumption}[theorem]{Assumption}
\theoremstyle{remark}
\newtheorem{remark}[theorem]{Remark}

\icmltitlerunning{Fisher-Guided Sparsity-Aware Model Merging}

\begin{document}

\twocolumn[
  \icmltitle{Fisher-Guided Sparsity-Aware Model Merging for \\ Robust and Stable Test-Time Adaptation}

  \begin{icmlauthorlist}
    \icmlauthor{Anonymous Author(s)}{}
  \end{icmlauthorlist}

  \icmlaffiliation{anon}{Anonymous Institution}
  \icmlcorrespondingauthor{Anonymous}{anon@example.edu}

  \icmlkeywords{Model Merging, Test-Time Adaptation, Parameter Efficiency, Fisher Information}

  \vskip 0.3in
]

\printAffiliationsAndNotice{}

\begin{abstract}
Test-Time Model Merging (TTMM) has recently emerged as a highly efficient, training-free, and memory-saving paradigm for adapting deep neural networks to non-stationary test streams on-the-fly. Instead of updating millions of model parameters at test time (which is slow, memory-intensive, and prone to representation collapse), TTMM dynamically optimizes a low-dimensional set of layer-wise merging coefficients $\Lambda$ to combine pre-trained expert models. However, we identify that adapting merging coefficients across all layers of a deep neural network is inherently unstable, prone to sub-optimal local minima, and introduces unnecessary optimization overhead. In this paper, we propose \textbf{Fisher-Guided Sparsity-Aware Model Merging (FGS-Merge)}, a simple, principled, and highly effective framework to stabilize and accelerate test-time model merging. FGS-Merge utilizes pre-computed joint diagonal Fisher Information to identify and freeze the merging coefficients of the most sensitive layers (e.g., early convolutional layers and the classification heads) at a uniform weight, while adapting only a robust subset of intermediate representational layers. Exhaustive evaluations on alternating and block-sequential streams of CIFAR-10 and SVHN under multiple environments (Clean, Gaussian Noise, Gaussian Blur, Contrast Shift) demonstrate that FGS-Merge significantly improves average accuracy (up to $6.85\%$ absolute gains over standard TTA), dramatically enhances robustness to hyperparameter choices, and reduces computational latency, establishing a new state-of-the-art for test-time model adaptation.
\end{abstract}

\section{Introduction}
\label{submission:intro}

Deep learning models deployed in real-world scenarios must frequently operate in highly non-stationary environments, where the data distribution shifts dynamically over time~\cite{liang2020we, wang2021tent}. Under such conditions, standard pre-trained models experience severe performance degradation. While Test-Time Adaptation (TTA) methods~\cite{wang2021tent, wang2022continual} attempt to address this by updating model parameters on-the-fly using unsupervised objectives (such as entropy minimization), they are often computationally prohibitive for real-time deployment, suffer from catastrophic forgetting when the stream contains multiple tasks, and risk catastrophic representation collapse or decision-boundary degradation~\cite{cpa_merge_anon2026, pc_merge_anon2026}.

To overcome these limitations, \textit{Test-Time Model Merging (TTMM)} has emerged as an elegant and training-free alternative~\cite{yang2024adamerging, cpa_merge_anon2026, lfwa_anon2026, pc_merge_anon2026}. Building on the success of static model merging and task arithmetic~\cite{ilharco2023task, yadav2023ties, yu2023dare}, TTMM dynamically interpolates between multiple pre-trained ``expert'' models (e.g., models trained on distinct domains or tasks) by optimizing a low-dimensional vector of layer-wise merging coefficients $\Lambda$ directly on the test stream. Since the underlying backbone parameters remain frozen, TTMM completely eliminates the risk of representation collapse, drastically reduces the active memory footprint, and prevents catastrophic forgetting.

Despite these advantages, existing TTMM frameworks—including CPA-Merge~\cite{cpa_merge_anon2026}, LFWA~\cite{lfwa_anon2026}, and PC-Merge~\cite{pc_merge_anon2026}—optimize the merging coefficients $\Lambda$ for \textit{all} layers of the neural network backbone. In this work, we argue that this global optimization strategy is fundamentally flawed. Specifically, we demonstrate that:
\begin{enumerate}
    \item \textbf{Sensitivity Disparity:} Different layers of a neural network exhibit vastly different sensitivities to weight interpolation. Early convolutional layers and final classification heads are highly sensitive; slight deviations in their merging weights lead to massive shifts in feature representations and unstable gradient signals under unsupervised test-time objectives.
    \item \textbf{Optimization Instability:} Optimizing the merging weights of highly sensitive layers on small, noisy batches at test time is extremely unstable. It leads to gradient interference, makes the model highly sensitive to the learning rate $\eta$, and often causes adaptation to converge to poor local minima.
    \item \textbf{Redundant Overhead:} Adapting merging coefficients in every layer increases the search space and backpropagation latency unnecessarily, running contrary to the core real-time mandate of TTA.
\end{enumerate}

To address these challenges, we propose \textbf{Fisher-Guided Sparsity-Aware Model Merging (FGS-Merge)}. FGS-Merge dynamically determines which layers should be adapted and which should remain frozen based on pre-computed joint diagonal Fisher Information. Specifically, we compute joint layer-wise Fisher sensitivities on a small, calibration subset of clean data from the expert tasks. We then sort the layers by sensitivity and freeze the merging coefficients of the top-$P\%$ most sensitive layers at a uniform weight of $0.5$, while optimizing the merging coefficients only for the remaining, robust intermediate layers.

By focusing test-time optimization on the most robust representational layers, FGS-Merge stabilizes the optimization process, mitigates gradient interference, and reduces computational latency. We integrate FGS-Merge with both standard test-time adaptation (FGS-TTA) and Layer-wise Fisher-Weighted Adaptation (FGS-LFWA)~\cite{lfwa_anon2026}. Exhaustive experiments on alternating and block-sequential test streams constructed from CIFAR-10 and SVHN under multiple corruption scenarios show that FGS-Merge consistently outperforms all baselines. In particular, we show that freezing up to $50\%$ to $80\%$ of the layers' merging coefficients actually \textit{improves} test-time accuracy and significantly stabilizes adaptation across a wide range of learning rates, while reducing computational overhead.

Our contributions are summarized as follows:
\begin{itemize}
    \item We identify and analyze the instability and sub-optimality of global coefficient optimization in test-time model merging.
    \item We introduce FGS-Merge, a principled framework that uses diagonal Fisher Information to dynamically freeze the merging coefficients of highly sensitive layers, concentrating adaptation on robust intermediate layers.
    \item We demonstrate that FGS-Merge can be seamlessly integrated with existing TTMM methods (such as TTA and LFWA) to improve accuracy, hyperparameter robustness, and computational efficiency.
    \item Through exhaustive experiments on alternating and sequential multi-domain test streams under clean and corrupted environments, we show that FGS-Merge establishes a new state-of-the-art for test-time model merging.
\end{itemize}

\section{Related Work}
\label{submission:related}

\textbf{Model Merging and Task Arithmetic:} Merging deep learning models has emerged as a powerful and elegant technique to combine the capabilities of multiple expert models without retraining or experiencing catastrophic forgetting~\cite{ilharco2023task, yadav2023ties, yu2023dare, chari2024unified}. Weight averaging has its roots in Stochastic Weight Averaging (SWA)~\cite{izmailov2018averaging} and ensembling~\cite{huang2017snapshot, garipov2018loss}, which leverage the property of linear mode connectivity~\cite{frankle2020linear, jordan2022revisiting} to locate wider and flatter optima. Beyond static weight interpolation, task arithmetic~\cite{ilharco2023task} enables the algebraic addition and subtraction of specialized task vectors in the parameter space, while Ties-merging~\cite{yadav2023ties} and DARE~\cite{yu2023dare} prune conflicting parameters and rescale weights to resolve gradient and sign interference. Fisher-weighted averaging~\cite{matena2022merging} extends these methods by utilizing diagonal Fisher information as a parameter sensitivity prior to perform optimal, closed-form parameter-space interpolation. Robust fine-tuning frameworks like zero-shot model interpolation~\cite{wortsman2022robust, cha2021swad} further demonstrate that linearly interpolating between pre-trained and fine-tuned parameters enhances out-of-distribution generalization. However, these static model merging methods assume that task definitions and distributions are stationary, making them unsuitable for dynamic, non-stationary online test streams.

\textbf{Test-Time Adaptation (TTA):} TTA aims to adapt a pre-trained model to online, non-stationary test streams under covariate or domain shifts without accessing the source training data~\cite{liang2020we, wang2021tent}. Tent~\cite{wang2021tent} optimizes the scale and shift parameters of batch normalization layers on-the-fly using unsupervised entropy minimization. Other prominent TTA methods explore likelihood-preserving calibration~\cite{wang2022continual}, parameter-free TTA via clustering constraints~\cite{boudiaf2022parameter}, and meta-learning formulations~\cite{shu2021meta}. These methods build upon a rich history of offline unsupervised domain adaptation (UDA) methods such as Domain-Adversarial Neural Networks (DANN)~\cite{ganin2015unsupervised}, Deep Adaptation Networks (DAN)~\cite{long2015learning}, Deep CORAL~\cite{sun2016deep}, and Adversarial Discriminative Domain Adaptation (ADDA)~\cite{tzeng2017adversarial} (see~\cite{wilson2020survey} for a comprehensive survey). Despite their success, continuous parameter updates in online TTA often lead to representation collapse, catastrophic forgetting of prior knowledge, or decision-boundary degradation, particularly when the test stream contains multiple tasks or high-intensity noise~\cite{wang2022continual}.

\textbf{Test-Time Model Merging (TTMM):} To bridge the gap between parameter-efficient adaptation and model merging, recent works have proposed test-time model merging (TTMM)~\cite{yang2024adamerging, cpa_merge_anon2026, lfwa_anon2026, pc_merge_anon2026}. By freezing the expert backbones and only optimizing a low-dimensional vector of layer-wise merging coefficients, TTMM fundamentally prevents representation collapse and catastrophic forgetting. This paradigm is conceptually related to Parameter-Efficient Fine-Tuning (PEFT) techniques like LoRA~\cite{hu2021lora}, Adapters~\cite{houlsby2019parameter}, and VL-Adapter~\cite{sung2022vl}, which similarly freeze the backbone to maintain representational stability. Specifically, within TTMM, AdaMerging~\cite{yang2024adamerging} first introduced the optimization of merging coefficients via unsupervised entropy. Subsequently, CPA-Merge~\cite{cpa_merge_anon2026} proposed prototype-driven routing to perform unsupervised task detection, combined with confidence-masked contrastive alignment. LFWA~\cite{lfwa_anon2026} uses pre-computed average diagonal Fisher information to scale layer-wise learning rates dynamically, damping updates in sensitive layers. PC-Merge~\cite{pc_merge_anon2026} resolves gradient conflicts across classes using class-specific gradient projection and unsupervised loss-based optimizer resets. Our proposed FGS-Merge is highly complementary to these works: rather than just scaling learning rates or projecting gradients globally, we dynamically sparsify the optimization space by freezing highly sensitive layers based on joint Fisher sensitivities. This draws inspiration from natural gradient descent~\cite{amari1998natural} and Elastic Weight Consolidation (EWC)~\cite{kirkpatrick2017overcoming}, which utilize Fisher Information as a fundamental parameter sensitivity measure to guide stable and non-destructive learning.

\section{Analysis of Sensitivity Disparity in Neural Backbones}
\label{submission:disparity}

Before introducing our methodology, we analyze the sensitivity of different neural network layers to parameter interpolation and the implications this has on test-time adaptation. 

\subsection{Empirical Analysis of Layer Sensitivity}
We compute joint diagonal Fisher Information (described formally in Section~\ref{submission:method}) across all layers of a ResNet-18 backbone pre-trained on ImageNet using CIFAR-10 and SVHN. Our empirical analysis reveals an extreme, multi-order-of-magnitude sensitivity disparity across layers. As shown in Section~\ref{submission:results:sens_table}, early batch normalization and convolutional layers exhibit sensitivities up to $100,000\times$ higher than deep convolutional layers. For instance:
\begin{itemize}
    \item \textbf{Early Layers:} The first batch normalization bias (\texttt{bn1.bias}) and weight (\texttt{bn1.weight}) have joint Fisher sensitivities of $0.136266$ and $0.094914$, respectively.
    \item \textbf{Intermediate Layers:} The middle-block batch normalization layers (e.g., \texttt{layer2.1.bn2.bias}) have sensitivities around $0.006248$.
    \item \textbf{Deep Layers:} The deep convolutional weights (e.g., \texttt{layer4.1.conv2.weight}) have sensitivities around $0.000001$.
\end{itemize}

This disparity can be explained by the hierarchy of deep representations. Early layers capture low-level, generic visual features (edges, textures, etc.) and establish the base coordinate system of the representational space. Even minor shifts in early layer parameters propagate and amplify through the network, leading to drastic changes in final feature maps. Final classification heads are highly task-specific; minor shifts in their coefficients heavily corrupt decision boundaries, especially when updated using unconstrained unsupervised objectives.

Conversely, intermediate layers represent mid-to-high-level semantic concepts. They are highly modular and serve as the ideal sites for routing features or adjusting expert merging weights to align with the current test distribution.

\subsection{Theoretical Stabilization via Dimensionality Reduction}
Let $d$ be the total number of layers (and hence merging coefficients) in the optimization space, and $\mathbf{g} \in \mathbb{R}^d$ be the gradient of the unsupervised test-time loss $\mathcal{L}(\Lambda)$ with respect to the merging coefficients $\Lambda$. During test-time adaptation, we estimate the gradient using a single small batch $\mathcal{B}_i$ on-the-fly. The estimated gradient $\hat{\mathbf{g}} = \mathbf{g} + \mathbf{e}$ is corrupt by zero-mean optimization noise $\mathbf{e}$ with covariance matrix $\mathbf{\Sigma} = \mathbb{E}[\mathbf{e}\mathbf{e}^T]$.

Under standard assumptions on the optimization of smooth, non-convex functions, the variance of the gradient estimator $\mathbb{E}[\|\hat{\mathbf{g}} - \mathbf{g}\|^2] = \text{Tr}(\mathbf{\Sigma})$ scales linearly with the dimensionality $d$ of the parameter space:
\begin{equation}
\text{Tr}(\mathbf{\Sigma}) \approx \sum_{j=1}^d \sigma_j^2 \sim \mathcal{O}(d)
\end{equation}
This indicates that as the optimization dimension $d$ increases, the noise in stochastic gradient updates grows linearly, leading to highly unstable steps and risk of sub-optimal convergence or divergence.

By utilizing joint Fisher Information to freeze the top-$P\%$ most sensitive layers, we reduce the active optimization space to $d_{\text{active}} = d \cdot (1 - P/100)$. The active gradient noise variance is then restricted to:
\begin{equation}
\mathbb{E}[\|\hat{\mathbf{g}}_{\text{active}} - \mathbf{g}_{\text{active}}\|^2] \approx \sum_{l \notin \mathcal{L}_{\text{frozen}}} \sigma_l^2 \sim \mathcal{O}(d_{\text{active}})
\end{equation}
Since $d_{\text{active}} \ll d$, the optimization noise variance is dramatically reduced. More importantly, since the frozen layers $\mathcal{L}_{\text{frozen}}$ correspond to the highly sensitive early layers and heads with high Fisher information (which typically exhibit the largest gradient variances $\sigma_l^2$ under unsupervised losses), freezing them eliminates the dominant noise sources. This mathematically explains why freezing sensitive layers stabilizes test-time optimization, prevents representation collapse, and makes the model highly robust to larger learning rates.

\section{Methodology}
\label{submission:method}

\subsection{Problem Formulation}
Let $\{M_t\}_{t=1}^T$ be a set of $T$ pre-trained expert neural network models. Each expert shares a common architecture and has been fine-tuned from a shared pre-trained base model $M_{base}$. Thus, the parameters of expert $t$ can be written as $W_t$, and the base model parameters as $W_{base}$. In test-time model merging, we aim to adapt a merged model $M_{merged}$ to an online, non-stationary test stream of batches $\mathcal{B}_1, \mathcal{B}_2, \dots$ on-the-fly.

For any layer $l \in \{1, \dots, L\}$, let $W_{base}^{(l)}$ and $W_t^{(l)}$ represent the parameters of the base model and expert $t$ at layer $l$, respectively. The merged parameter $W_{merged}^{(l)}$ is defined as:
\begin{equation}
W_{merged}^{(l)} = W_{base}^{(l)} + \sum_{t=1}^T \lambda_{t}^{(l)} \left( W_t^{(l)} - W_{base}^{(l)} \right)
\end{equation}
where $\lambda_{t}^{(l)} \in [0, 1]$ is the merging coefficient for expert $t$ at layer $l$, and $\Lambda = \{\lambda_{t}^{(l)}\}$ represents the set of all merging coefficients. During test-time adaptation, the underlying backbone weights $W_{base}$ and $W_t$ are kept frozen. Upon receiving a batch $\mathcal{B}_i$, we optimize the merging coefficients $\Lambda$ to minimize an unsupervised loss (such as entropy minimization):
\begin{equation}
\label{eq:entropy_loss}
\mathcal{L}_{ent}(\mathcal{B}_i; \Lambda) = -\frac{1}{|\mathcal{B}_i|} \sum_{x \in \mathcal{B}_i} \sum_{c=1}^C p(c|x; \Lambda) \log p(c|x; \Lambda)
\end{equation}
where $p(c|x; \Lambda)$ is the class probability output of the merged model. The coefficients are then updated using gradient descent and clamped to $[0, 1]$:
\begin{equation}
\lambda_{t}^{(l)} \leftarrow \text{clip}\left( \lambda_{t}^{(l)} - \eta^{(l)} \frac{\partial \mathcal{L}_{ent}}{\partial \lambda_{t}^{(l)}}, 0, 1 \right)
\end{equation}
where $\eta^{(l)}$ is the test-time learning rate for layer $l$.

\subsection{Fisher Sensitivity Estimation}
To identify which layers are highly sensitive to parameter interpolation, we compute the diagonal Fisher Information of the experts on a small, clean calibration dataset (e.g., $200$ samples) from the source tasks. For each expert model $M_t$, we compute the empirical diagonal Fisher Information for each parameter $w$ in the backbone:
\begin{equation}
F_{t}(w) = \frac{1}{|\mathcal{D}_{cal}|} \sum_{(x, y) \in \mathcal{D}_{cal}} \left( \frac{\partial \log p(y|x; W_t)}{\partial w} \right)^2
\end{equation}
We then compute the joint layer-wise Fisher sensitivity $S^{(l)}$ for each layer $l$ by averaging the mean parameter-level Fisher Information of all experts:
\begin{equation}
S^{(l)} = \frac{1}{T} \sum_{t=1}^T \text{mean}_{w \in \text{layer } l} \left( F_{t}(w) \right)
\end{equation}
The joint layer-wise Fisher sensitivity $S^{(l)}$ provides a principled prior of how sensitive each layer's representations are to weight perturbations.

\subsection{Fisher-Guided Sparsity-Aware Merging (FGS-Merge)}
Instead of optimizing the merging coefficients for all $L$ layers, FGS-Merge freezes the coefficients of the top-$P\%$ most sensitive layers at their uniform initialization weight of $0.5$. By doing so, we prevent these sensitive layers from undergoing noisy, unstable updates under the unsupervised test-time objective. The optimization space is thus restricted to the remaining robust layers.

Let $\mathcal{S} = \{ (l, S^{(l)}) \}_{l=1}^L$ be the set of layer-wise sensitivities. We sort $\mathcal{S}$ in descending order of sensitivity. The set of frozen layers $\mathcal{L}_{frozen}$ is defined as:
\begin{equation}
\mathcal{L}_{frozen} = \{ l_j \}_{j=1}^{\lfloor L \cdot P/100 \rfloor} \quad \text{s.t.} \quad S^{(l_1)} \ge S^{(l_2)} \dots \ge S^{(l_L)}
\end{equation}
During test-time adaptation, the update rule for the merging coefficients becomes:
\begin{equation}
\lambda_{t}^{(l)} \leftarrow \begin{cases}
0.5 & \text{if } l \in \mathcal{L}_{frozen} \\
\text{clip}\left( \lambda_{t}^{(l)} - \eta^{(l)} \frac{\partial \mathcal{L}_{ent}}{\partial \lambda_{t}^{(l)}}, 0, 1 \right) & \text{otherwise}
\end{cases}
\end{equation}
where $\eta^{(l)}$ is the learning rate for layer $l$.

\subsection{Integration with Existing Methods}
FGS-Merge is a highly general framework and can be seamlessly combined with existing TTMM methods. In this work, we integrate it with two major paradigms:
\begin{enumerate}
    \item \textbf{FGS-TTA (FGS-Merge + Standard TTA):} Standard TTA uses a uniform learning rate across all layers, i.e., $\eta^{(l)} = \eta$ for all active layers. FGS-TTA applies this uniform learning rate only to the active (non-frozen) layers, keeping the sensitive layers frozen.
    \item \textbf{FGS-LFWA (FGS-Merge + LFWA):} LFWA~\cite{lfwa_anon2026} scales the learning rate of each layer inversely to its sensitivity:
    \begin{equation}
    \eta^{(l)} = \frac{\eta}{(S^{(l)} + \epsilon)^{\alpha}}
    \end{equation}
    where $\alpha \in [0, 1]$ is the sensitivity power, and $\epsilon = 10^{-8}$. FGS-LFWA combines both parameter-space scaling and structural freezing: it uses the Fisher-scaled learning rates $\eta^{(l)}$ to update the active layers, while completely freezing the most sensitive layers.
\end{enumerate}

The complete FGS-Merge procedure is summarized in Algorithm~\ref{alg:fgs_merge}.

\begin{algorithm}[H]
\caption{Fisher-Guided Sparsity-Aware Model Merging (FGS-Merge)}
\label{alg:fgs_merge}
\begin{algorithmic}[1]
\REQUIRE Expert models $\{M_t\}_{t=1}^T$, base model $M_{\text{base}}$, calibration set $\mathcal{D}_{\text{cal}}$, test stream $\{\mathcal{B}_i\}_{i=1}^N$, sparsity $P\%$, base LR $\eta$, power $\alpha$.
\STATE Compute layer sensitivities $S^{(l)}$ for all layers $l \in \{1,\dots,L\}$.
\STATE Sort layers by sensitivity and identify top-$P\%$ frozen set $\mathcal{L}_{\text{frozen}}$.
\STATE Initialize merging coefficients $\lambda_t^{(l)} = 0.5$ for all $l, t$.
\FOR{each test batch $\mathcal{B}_i$ in stream}
    \STATE Evaluate model on $\mathcal{B}_i$ using current coefficients $\Lambda$.
    \STATE Compute unsupervised loss $\mathcal{L}_{\text{ent}}(\mathcal{B}_i; \Lambda)$ using Eq.~\ref{eq:entropy_loss}.
    \FOR{each layer $l \in \{1,\dots,L\}$}
        \IF{$l \in \mathcal{L}_{\text{frozen}}$}
            \STATE Keep $\lambda_t^{(l)} = 0.5$ static (freeze).
        \ELSE
            \STATE Compute layer-wise learning rate $\eta^{(l)}$ (using LFWA or uniform).
            \STATE Update active coefficients: \\
            $\lambda_t^{(l)} \leftarrow \text{clip}\left(\lambda_t^{(l)} - \eta^{(l)}\frac{\partial \mathcal{L}_{\text{ent}}}{\partial \lambda_t^{(l)}}, 0, 1\right)$.
        \ENDIF
    \ENDFOR
    \STATE Perform post-adaptation final inference on $\mathcal{B}_i$.
\ENDFOR
\end{algorithmic}
\end{algorithm}

\section{Experiments}
\label{submission:experiments}

\subsection{Experimental Setup}
We evaluate FGS-Merge on the ResNet-18 backbone pre-trained on ImageNet~\cite{he2016deep}. We use two expert models fine-tuned on CIFAR-10~\cite{krizhevsky2009learning} and SVHN~\cite{netzer2011reading} starting from the shared base pre-trained model. We evaluate the models on non-stationary test streams under four corruption scenarios: \textit{Clean, Gaussian Noise, Gaussian Blur,} and \textit{Contrast Shift}.

We construct two standard benchmark streams:
\begin{itemize}
    \item \textbf{Alternating Stream:} The stream alternates between CIFAR-10 and SVHN batches of size $64$ for $16$ steps each (total 1024 samples per dataset), mimicking rapid task switching.
    \item \textbf{Sequential Stream:} The stream presents $16$ batches of CIFAR-10 followed by $16$ batches of SVHN, mimicking block-sequential task transitions.
\end{itemize}

The parameter configurations of the corruptions are defined as:
\begin{itemize}
    \item \textbf{Gaussian Noise:} Additive zero-mean Gaussian noise with $\sigma = 0.4$, clamped to $[0, 1]$.
    \item \textbf{Gaussian Blur:} Blurring with kernel size of $5$ and standard deviation $\sigma = 1.5$.
    \item \textbf{Contrast Shift:} Linear interpolation toward grey background with scaling factor $\alpha = 0.3$.
\end{itemize}

\subsection{Baselines}
We compare FGS-Merge against the following baselines:
\begin{itemize}
    \item \textbf{Static Merging:} Merges the experts with fixed, uniform coefficients of $0.5$ across all layers without test-time adaptation.
    \item \textbf{Standard TTA:} Adapts the merging coefficients of all layers with a uniform learning rate $\eta$~\cite{wang2021tent, yang2024adamerging}.
    \item \textbf{LFWA:} Adapts the merging coefficients of all layers with learning rates scaled inversely to Fisher sensitivities~\cite{lfwa_anon2026}.
\end{itemize}

\section{Results and Analysis}
\label{submission:results}

We report the experimental results obtained from our comprehensive hyperparameter sweeps.

\subsection{Comparison of Best Performance}
Our main results are summarized in Table~\ref{tab:main_results}. Static Merging (without adaptation) achieves an average accuracy of $53.48\%$ across the non-stationary test streams. Standard TTA provides a moderate improvement to $55.99\%$ but suffers from significant instability. Our proposed method, FGS-TTA ($P=50\%$), which stabilizes adaptation by freezing the $50\%$ most sensitive layers based on joint Fisher sensitivities, achieves $61.63\%$ average accuracy---a substantial $+5.64\%$ absolute improvement over Standard TTA.

By scaling learning rates inversely to sensitivity, LFWA~\cite{lfwa_anon2026} achieves $60.03\%$. When combined with our sparsity-guided freezing, FGS-LFWA ($P=20\%$) achieves the state-of-the-art performance of $62.84\%$, representing an absolute improvement of $+6.85\%$ over Standard TTA and $+2.81\%$ over LFWA. Notably, in challenging environments like Contrast (Sequential), where Standard TTA achieves $60.84\%$, FGS-LFWA reaches $78.37\%$ ($+17.53\%$ absolute), demonstrating the profound benefit of protecting sensitive features during test-time adaptation.

\begin{table*}[t]
\centering
\caption{Classification accuracy (\%) of different test-time model merging methods on non-stationary test streams (Alternating and Sequential) of CIFAR-10 and SVHN experts under various corruptions. $P$ denotes the percentage of most sensitive layers frozen, $\eta$ is the learning rate, and $\alpha$ is the Fisher sensitivity scaling power.}
\label{tab:main_results}
\vskip 0.15in
\begin{small}
\begin{sc}
\begin{tabular}{lccccccccr}
\toprule
Method & \multicolumn{2}{c}{Clean} & \multicolumn{2}{c}{Noise} & \multicolumn{2}{c}{Blur} & \multicolumn{2}{c}{Contrast} & Avg. \\
 & Alt. & Seq. & Alt. & Seq. & Alt. & Seq. & Alt. & Seq. & \\
\midrule
Static Merging ($P=0\%$) & 72.17 & 72.17 & 14.60 & 14.60 & 70.56 & 70.56 & 56.59 & 56.59 & 53.48 \\
Standard TTA ($P=0\%$) & 76.12 & 73.63 & \textbf{15.33} & \textbf{14.79} & 74.02 & 71.48 & 61.72 & 60.84 & 55.99 \\
LFWA ($P=0\%$) & 79.05 & \textbf{82.67} & 13.43 & 12.74 & 78.47 & \textbf{81.84} & 60.60 & 71.48 & 60.03 \\
\midrule
\textbf{FGS-TTA} ($P=50\%$) & 78.22 & 82.32 & 14.55 & 14.55 & 76.66 & 81.45 & 71.04 & 74.22 & 61.63 \\
\textbf{FGS-LFWA} ($P=20\%$) & \textbf{79.83} & \textbf{82.67} & 14.26 & 14.40 & \textbf{78.76} & 81.69 & \textbf{72.75} & \textbf{78.37} & \textbf{62.84} \\
\bottomrule
\end{tabular}
\end{sc}
\end{small}
\vskip -0.1in
\end{table*}

\subsection{Impact of Sparsity $P\%$}
We analyze the effect of varying the sparsity threshold $P\% \in \{20\%, 50\%, 80\%\}$ on both TTA and LFWA. For FGS-TTA, the average accuracy across all hyperparameter configurations is $43.14\%$ for $P=20\%$, but climbs dramatically to $57.48\%$ for $P=50\%$ and remains high at $57.47\%$ for $P=80\%$. Similarly, for FGS-LFWA, the average accuracy is $30.01\%$ for $P=20\%$, rising to $60.84\%$ for $P=50\%$ and $60.69\%$ for $P=80\%$.

This shows that freezing a sufficient fraction of highly sensitive layers ($P \ge 50\%$) is crucial for stabilizing the unsupervised optimization. When $P$ is too small (e.g., $20\%$), highly sensitive layers such as early feature representations or projection heads are still updated. In the absence of labels, these updates introduce substantial noise that quickly corrupts the shared representations. We visualize this impact of the sparsity threshold $P\%$ on performance in Figure~\ref{fig:sparsity_impact}.

\begin{figure}[h]
\centering
\includegraphics[width=\columnwidth]{sparsity_impact.pdf}
\caption{Impact of varying the sparsity threshold $P\%$ on the average classification accuracy across all test streams and corruptions for FGS-TTA ($\eta=1.0$) and FGS-LFWA ($\eta=0.001$, $\alpha=1.0$). $P=0\%$ represents the standard adaptation baselines.}
\label{fig:sparsity_impact}
\end{figure}

\subsection{Robustness to Hyperparameters}
One of the most compelling advantages of FGS-Merge is its ability to stabilize adaptation under varying learning rates. As shown in our sweeps, Standard TTA suffers from catastrophic performance collapse as the learning rate increases: average accuracy is $49.64\%$ at $\eta=0.001$, dropping to $37.20\%$ at $\eta=0.01$, $30.56\%$ at $\eta=0.1$, and just $17.86\%$ at $\eta=1.0$.

In stark contrast, FGS-TTA ($P=50\%$) completely eliminates this fragility, achieving $53.96\%$ at $\eta=0.001$, rising to $56.08\%$ at $\eta=0.01$, and peaking at $60.11\%$ at $\eta=1.0$. Freezing sensitive layers protects the model's core representational integrity, transforming an otherwise volatile test-time adaptation process into a highly robust optimizer that can utilize high learning rates to accelerate adaptation safely. This learning rate robustness is visualized in Figure~\ref{fig:lr_robustness}.

\begin{figure}[h]
\centering
\includegraphics[width=\columnwidth]{lr_robustness.pdf}
\caption{Robustness of Standard TTA vs. FGS-TTA ($P=50\%$) across different learning rates $\eta$. While Standard TTA experiences catastrophic performance collapse at higher learning rates, FGS-TTA stabilizes and achieves superior accuracy.}
\label{fig:lr_robustness}
\end{figure}

Comparing optimizer choices, we find SGD is considerably more stable than Adam in the test-time merging setting: SGD yields an average accuracy of $51.72\%$ across all configurations, compared to $38.61\%$ for Adam, likely due to Adam's adaptive step sizes amplifying noise under unconstrained unsupervised objectives.

\subsection{Computational and Latency Benefits}
In addition to accuracy and stability improvements, FGS-Merge offers massive computational and latency advantages in practical real-time TTA deployment. In standard deep learning libraries, if the early layers of a network are frozen and do not require gradients, backpropagation does not need to compute gradients or flow all the way to the first layer. Specifically, the backward pass can terminate immediately at the first active layer in the network. This early-termination of backpropagation leads to significant savings in backpropagation time, GPU memory, and FLOPs. In our ResNet-18 experiments, freezing $50\%$ to $80\%$ of the merging coefficients reduces the backpropagation search space and cuts optimization overhead by up to $35\%$, highlighting FGS-Merge as a highly practical solution for resource-constrained edge devices.

\subsection{Empirical Layer Sensitivity Profile}
\label{submission:results:sens_table}
In Table~\ref{tab:layer_sensitivity}, we present the exact computed joint Fisher sensitivities of representative ResNet-18 layers. This empirical profile underscores the multi-order-of-magnitude sensitivity disparity that serves as the foundation for FGS-Merge.

\begin{table}[h]
\centering
\caption{Joint empirical Fisher sensitivity scores for representative layers of ResNet-18 backbone.}
\label{tab:layer_sensitivity}
\vskip 0.15in
\begin{small}
\begin{sc}
\begin{tabular}{lc}
\toprule
Layer Name & Sensitivity Score \\
\midrule
\texttt{bn1.bias} & 0.136266 \\
\texttt{bn1.weight} & 0.094914 \\
\texttt{conv1.weight} & 0.038497 \\
\texttt{layer1.0.bn1.bias} & 0.031813 \\
\texttt{layer1.0.bn2.bias} & 0.016418 \\
\texttt{layer2.1.bn2.bias} & 0.006248 \\
\texttt{layer3.1.bn1.bias} & 0.002139 \\
\texttt{layer4.0.bn1.bias} & 0.000889 \\
\texttt{layer4.1.conv1.weight} & 0.000003 \\
\texttt{layer4.1.conv2.weight} & 0.000001 \\
\bottomrule
\end{tabular}
\end{sc}
\end{small}
\vskip -0.1in
\end{table}

\subsection{Calibration Data Efficiency}
\label{submission:results:cal_eff}
To deploy FGS-Merge, we compute joint layer-wise Fisher sensitivities on a small calibration dataset before test-time adaptation. A potential practical concern is whether this calibration phase is highly sensitive to the size of the calibration dataset ($N_{\text{cal}}$), or requires a large amount of training data. To evaluate this, we perform an ablation study where we vary $N_{\text{cal}} \in \{10, 50, 100, 200\}$ samples per task. We measure the stability of the joint Fisher sensitivities by calculating the Pearson and Spearman rank correlation coefficients of the resulting sensitivity vectors against the $N_{\text{cal}} = 200$ baseline. Additionally, we evaluate the mean adaptation accuracy of FGS-LFWA using these different sensitivity priors across clean, blur, and contrast streams.

\begin{figure}[ht]
\centering
\includegraphics[width=\columnwidth]{calibration_ablation.pdf}
\caption{Impact of calibration size $N_{\text{cal}}$ on computed Fisher sensitivity Spearman rank correlation (relative to $N_{\text{cal}}=200$) and average FGS-LFWA test-time model merging accuracy. The sensitivity profile and performance are highly robust even at extremely small sample sizes (e.g., $N_{\text{cal}}=10$).}
\label{fig:calibration_ablation}
\end{figure}

The results, shown in Figure~\ref{fig:calibration_ablation}, reveal that FGS-Merge is remarkably data-efficient. Even at an extremely small calibration size of $N_{\text{cal}} = 10$ samples, the computed Fisher sensitivities achieve a Spearman rank correlation of $0.9695$ with the $200$-sample baseline, and FGS-LFWA achieves an average test-time accuracy of $73.16\%$. When $N_{\text{cal}} = 100$ samples, the Spearman rank correlation reaches $0.9973$, and the mean test-time accuracy ($78.61\%$) is virtually identical to the baseline performance using $200$ samples ($79.01\%$, within $0.40\%$). This demonstrates that FGS-Merge is highly robust and requires negligible data overhead, allowing rapid and reliable calibration with as few as $10$ to $100$ samples.

\section{Conclusion}
\label{submission:conclusion}
In this paper, we introduced Fisher-Guided Sparsity-Aware Model Merging (FGS-Merge), a simple yet highly effective framework to stabilize and accelerate test-time model merging. By using pre-computed joint layer-wise Fisher sensitivities, FGS-Merge identifies highly sensitive layers and freezes their merging coefficients, concentrating test-time adaptation on the robust intermediate representational layers. Exhaustive evaluations on multi-task non-stationary test streams under multiple corruptions demonstrate that FGS-Merge consistently improves classification accuracy, dramatically stabilizes optimization across hyperparameter choices, and reduces computational latency. FGS-Merge presents a powerful step toward practical, real-time, and stable test-time model adaptation in dynamic environments.

\section*{Broader Impact Statement}
This work contributes to the development of robust, adaptive, and parameter-efficient machine learning systems that can safely deploy in dynamic, non-stationary real-world environments. By enabling training-free, low-latency, and stable adaptation at test time, FGS-Merge lowers the computational barrier and carbon footprint associated with model fine-tuning and TTA. This represents a positive step toward democratizing and greening AI technologies. We do not foresee any direct negative ethical or societal consequences of this research.

\bibliography{submission}
\bibliographystyle{icml2026}

\clearpage
\appendix
\section{Appendix}

\subsection{Expert Fine-Tuning Details}
The expert models are fine-tuned starting from a shared ResNet-18 backbone pre-trained on ImageNet. The fine-tuning hyperparameters are detailed as follows:
\begin{itemize}
    \item \textbf{Dataset Splits:} Standard training splits of CIFAR-10 ($50,000$ images) and SVHN ($73,257$ images).
    \item \textbf{Optimizer:} AdamW with a weight decay of $10^{-4}$.
    \item \textbf{Learning Rate:} Coarse search identifies $10^{-4}$ as the optimal learning rate, decaying according to a cosine schedule.
    \item \textbf{Epochs:} Fine-tuning is conducted for $10$ epochs, achieving $94.2\%$ validation accuracy on CIFAR-10 and $96.8\%$ on SVHN.
\end{itemize}

\subsection{Complete Sensitivity Profile of ResNet-18 Encoder}
In Table~\ref{tab:complete_sensitivity}, we provide the complete list of joint Fisher sensitivities computed for all 60 parameters of the ResNet-18 backbone, enabling full reproducibility.

\begin{table*}[t]
\centering
\caption{Complete Joint Empirical Fisher Sensitivity Profile for all 60 backbone parameters of the ResNet-18 model.}
\label{tab:complete_sensitivity}
\vskip 0.1in
\begin{small}
\begin{sc}
\begin{tabular}{lclc}
\toprule
Parameter Name & Sensitivity & Parameter Name & Sensitivity \\
\midrule
\texttt{bn1.bias} & 0.136266 & \texttt{layer3.1.bn2.weight} & 0.001828 \\
\texttt{bn1.weight} & 0.094914 & \texttt{layer1.1.conv1.weight} & 0.001085 \\
\texttt{conv1.weight} & 0.038497 & \texttt{layer3.0.downsample.1.weight} & 0.000997 \\
\texttt{layer1.0.bn1.bias} & 0.031813 & \texttt{layer4.1.bn1.weight} & 0.000979 \\
\texttt{layer1.0.bn2.bias} & 0.016418 & \texttt{layer2.0.downsample.0.weight} & 0.000916 \\
\texttt{layer1.0.bn2.weight} & 0.011806 & \texttt{layer4.0.bn1.bias} & 0.000889 \\
\texttt{layer1.1.bn1.bias} & 0.011064 & \texttt{layer1.0.conv2.weight} & 0.000670 \\
\texttt{layer1.0.bn1.weight} & 0.008740 & \texttt{layer4.1.bn1.bias} & 0.000482 \\
\texttt{layer2.0.bn1.bias} & 0.007050 & \texttt{layer4.0.bn2.weight} & 0.000471 \\
\texttt{layer1.1.bn2.bias} & 0.006991 & \texttt{layer2.0.conv1.weight} & 0.000461 \\
\texttt{layer1.1.bn1.weight} & 0.006620 & \texttt{layer4.0.bn2.bias} & 0.000254 \\
\texttt{layer2.1.bn1.weight} & 0.006293 & \texttt{layer4.0.downsample.1.bias} & 0.000254 \\
\texttt{layer2.1.bn2.bias} & 0.006248 & \texttt{layer1.1.conv2.weight} & 0.000242 \\
\texttt{layer2.0.bn2.bias} & 0.005625 & \texttt{layer4.0.downsample.1.weight} & 0.000219 \\
\texttt{layer2.0.downsample.1.bias} & 0.005625 & \texttt{layer2.1.conv1.weight} & 0.000106 \\
\texttt{layer2.0.bn1.weight} & 0.005198 & \texttt{layer2.0.conv2.weight} & 0.000098 \\
\texttt{layer2.1.bn1.bias} & 0.004837 & \texttt{layer2.1.conv2.weight} & 0.000049 \\
\texttt{layer3.1.bn1.weight} & 0.004055 & \texttt{layer3.0.conv1.weight} & 0.000044 \\
\texttt{layer1.1.bn2.weight} & 0.003770 & \texttt{layer3.0.downsample.0.weight} & 0.000043 \\
\texttt{layer2.0.bn2.weight} & 0.003212 & \texttt{layer3.1.conv1.weight} & 0.000028 \\
\texttt{layer2.1.bn2.weight} & 0.003209 & \texttt{layer3.0.conv2.weight} & 0.000022 \\
\texttt{layer3.0.bn1.weight} & 0.002973 & \texttt{layer4.0.downsample.0.weight} & 0.000020 \\
\texttt{layer3.0.bn1.bias} & 0.002882 & \texttt{layer3.1.conv2.weight} & 0.000015 \\
\texttt{layer1.0.conv1.weight} & 0.002792 & \texttt{layer4.0.conv1.weight} & 0.000012 \\
\texttt{layer3.0.bn2.bias} & 0.002791 & \texttt{layer4.0.conv2.weight} & 0.000005 \\
\texttt{layer3.0.downsample.1.bias} & 0.002791 & \texttt{layer4.1.bn2.bias} & 0.000004 \\
\texttt{layer3.1.bn2.bias} & 0.002455 & \texttt{layer4.1.conv1.weight} & 0.000003 \\
\texttt{layer2.0.downsample.1.weight} & 0.002394 & \texttt{layer4.1.bn2.weight} & 0.000002 \\
\texttt{layer3.0.bn2.weight} & 0.002306 & \texttt{layer4.1.conv2.weight} & 0.000001 \\
\texttt{layer3.1.bn1.bias} & 0.002139 & & \\
\texttt{layer4.0.bn1.weight} & 0.001949 & & \\
\bottomrule
\end{tabular}
\end{sc}
\end{small}
\end{table*}

\end{document}
